% ==============================================================================
% Corpo Reutilizável do Cheat Sheet de Programação Dinâmica
% Compartilhado entre cheatsheet/cheatsheet.tex e relatorio/secoes/apendice_cheatsheet.tex
% Responsável: Integrante P5 (com tabela alimentada por P2)
% ==============================================================================

\section*{Guia Rápido: Programação Dinâmica em Uma Página}

% ------------------------------------------------------------------------------
% Bloco (a): Definição em Uma Frase
% ------------------------------------------------------------------------------
\subsection*{1. Definição Essencial}
% TODO(P5): Inserir a definição canônica concisa em uma única frase.
% Exemplo: "Técnica de projeto para problemas de otimização que resolve cada subproblema
% apenas uma vez, armazenando seu resultado em tabela para reaproveitamento futuro."
\noindent\fbox{%
  \parbox{\dimexpr\linewidth-2\fboxsep-2\fboxrule}{%
    \textbf{Definição:} \textit{[TODO(P5): Inserir a definição em uma frase aqui]}
  }%
}

\vspace{0.5cm}

% ------------------------------------------------------------------------------
% Bloco (b): Diagrama de Decisão "Quando utilizar esta abordagem?"
% ------------------------------------------------------------------------------
\subsection*{2. Diagrama de Decisão: Quando Utilizar?}
% TODO(P5): Completar o fluxograma de decisão em TikZ:
% 1. O problema é de otimização? (Não -> Buscar outra técnica; Sim -> Próximo)
% 2. Possui Subestrutura Ótima? (Não -> PD não se aplica; Sim -> Próximo)
% 3. Os subproblemas se sobrepõem?
%    - Se Não (independentes) -> Divisão e Conquista.
%    - Se Sim -> Uma escolha local gulosa é sempre globalmente ótima?
%      * Se Sim -> Algoritmo Guloso.
%      * Se Não -> PROGRAMAÇÃO DINÂMICA.

\begin{center}
\begin{tikzpicture}[
    node distance=1.5cm,
    decision/.style={diamond, draw, fill=blue!10, text width=5em, text centered, inner sep=0pt, font=\footnotesize},
    block/.style={rectangle, draw, fill=gray!10, text width=6em, text centered, rounded corners, minimum height=2em, font=\footnotesize},
    line/.style={draw, -Latex, font=\scriptsize}
]
  % Nós em esqueleto
  \node [decision] (opt) {Problema de Otimização?};
  \node [decision, right=2cm of opt] (sub) {Subestrutura Ótima?};
  \node [decision, below=1.2cm of sub] (over) {Subproblemas Sobrepostos?};
  \node [block, right=1.5cm of over] (pd) {\textbf{Programação Dinâmica}};
  \node [block, below=1cm of over] (dc) {Divisão e Conquista};
  \node [block, below=1.2cm of opt] (other) {Outras Abordagens};

  % Arestas conceituais
  \path [line] (opt) -- node[above]{Sim} (sub);
  \path [line] (opt) -- node[left]{Não} (other);
  \path [line] (sub) -- node[right]{Sim} (over);
  \path [line] (over) -- node[above]{Sim} (pd);
  \path [line] (over) -- node[right]{Não} (dc);
\end{tikzpicture}
\end{center}

\vspace{0.3cm}

% ------------------------------------------------------------------------------
% Bloco (c): Tabela Oficial de Complexidades
% ------------------------------------------------------------------------------
\subsection*{3. Tabela Comparativa de Complexidades}
% TODO(P5/P2): Preencher a tabela com os valores consolidados da Seção 4.
\begin{table}[h!]
\centering
\small
\begin{tabular}{@{}lllll@{}}
\toprule
\textbf{Problema} & \textbf{Ingênuo (Tempo)} & \textbf{PD (Tempo)} & \textbf{PD (Espaço)} & \textbf{Recorrência Chave} \\ \midrule
Fibonacci & $\Theta(2^n)$ & $\Theta(n)$ & $\Theta(n)$ ou $\Theta(1)$ & $F(n) = F(n-1) + F(n-2)$ \\
Corte de Hastes & $\Theta(2^n)$ & $\Theta(n^2)$ & $\Theta(n)$ & $r_n = \max (p_i + r_{n-i})$ \\
Cadeia de Matrizes & $\Omega(4^n / n^{3/2})$ & $\Theta(n^3)$ & $\Theta(n^2)$ & $m[i,j] = \min (m[i,k]+m[k+1,j]+p_{i-1}p_kp_j)$ \\
LCS & $O(2^{m+n})$ & $\Theta(m \cdot n)$ & $\Theta(m \cdot n)$ ou $\Theta(\min)$ & $c[i,j] = c[i-1,j-1]+1 \text{ ou } \max$ \\
Mochila 0/1 (Opcional) & $O(2^n)$ & $O(n \cdot W)$ & $O(n \cdot W)$ ou $O(W)$ & $V[i,w] = \max(V_{sem}, V_{com} + v_i)$ \\ \bottomrule
\end{tabular}
\end{table}

\vspace{0.3cm}

% ------------------------------------------------------------------------------
% Bloco (d): Esqueleto de Pseudocódigo (Top-Down vs. Bottom-Up)
% ------------------------------------------------------------------------------
\subsection*{4. Padrões de Implementação: Top-Down vs. Bottom-Up}
% TODO(P5): Inserir o template genérico dos dois estilos clássicos de PD.

\noindent
\begin{minipage}[t]{0.48\textwidth}
\textbf{Top-Down (com Memoização):}
\begin{algorithm}[H]
\caption{Padrao-Top-Down(x, memo)}
\SetKwFunction{FPD}{Padrao-Top-Down}
\If{caso base atingido}{
  \Return valor base\;
}
\If{x já está em memo}{
  \Return memo[x]\;
}
res $\leftarrow$ combinar escolhas recursivas\;
memo[x] $\leftarrow$ res\;
\Return res\;
\end{algorithm}
\end{minipage}
\hfill
\begin{minipage}[t]{0.48\textwidth}
\textbf{Bottom-Up (Tabulação Iterativa):}
\begin{algorithm}[H]
\caption{Padrao-Bottom-Up(n)}
Alocar tabela T[0..n]\;
Inicializar casos base em T\;
\For{tamanho $\leftarrow 1$ \KwTo $n$}{
  T[tamanho] $\leftarrow \min/\max$ sobre subproblemas já calculados\;
}
\Return T[n]\;
\end{algorithm}
\end{minipage}
